% Generated draft: requires mathematical and source review.
\documentclass[11pt]{article}
\usepackage[margin=25mm]{geometry}
\usepackage{fontspec}
\setmainfont{Latin Modern Roman}
\usepackage{amsmath,amssymb,mathtools,unicode-math}
\setmathfont{Latin Modern Math}
\usepackage{enumitem,needspace,xurl,hyperref}
\setlength{\parindent}{0pt}
\setlength{\parskip}{5pt}
\setlength{\emergencystretch}{3em}
\newcommand{\N}{\mathbb{N}}
\newcommand{\Z}{\mathbb{Z}}
\newcommand{\Q}{\mathbb{Q}}
\newcommand{\R}{\mathbb{R}}
\newcommand{\C}{\mathbb{C}}
\newcommand{\F}{\mathbb{F}}
\newcommand{\A}{\mathbb{A}}
\newcommand{\PP}{\mathbb{P}}
\newcommand{\E}{\mathbb{E}}
\newcommand{\eps}{\varepsilon}
\newcommand{\dd}{\,\mathrm{d}}
\begin{document}
\section*{UnsolvedMath ID: 800712}
\textbf{Source problem code:} catalogue entry
\subsection{Definitions and mathematical statement}
Let \(\ell\neq p\) be a prime and write
\[
\overline X=X\times_k k_{\mathrm{sep}},\qquad
V_\ell=\mathrm H^2(\overline X,\Q_\ell(1)),
\]
where \(k_{\mathrm{sep}}\) denotes a separable closure of \(k\), and \(\Q_\ell(1)\) denotes the first Tate twist. The supplied record refers to surfaces, but does not specify the hypotheses on \(X\) or \(k\), define \(p\) and \(q\), or explicitly identify \(k\) as a finite field of cardinality \(q\) and characteristic \(p\). These omissions must be resolved before this becomes a fully specified conjecture.

Following the supplied notation, let \(\operatorname{Pic}(X)\) be the group of line bundles on \(X\), let \(\operatorname{Pic}^0(X)\) denote its subgroup of classes algebraically equivalent to zero, and set
\[
\operatorname{NS}(X)=\operatorname{Pic}(X)/\operatorname{Pic}^0(X),\qquad
\rho(X)=\operatorname{rank}_{\Z}\operatorname{NS}(X).
\]
The record uses \(X\), rather than \(\overline X\), in this definition; no replacement by the geometric N\'eron--Severi group is made here.

Let \(F\) denote the action described in the record as geometric Frobenius on \(V_\ell\). The record calls \(P_2(X,t)\) its characteristic polynomial but supplies no determinant formula. In particular, the two polynomial conventions
\[
\det(t\,\mathrm{id}_{V_\ell}-F)
\quad\text{and}\quad
\det(\mathrm{id}_{V_\ell}-tF)
\]
have different root conventions and must be distinguished.

\textbf{Recorded assertion (Tate's Conjecture T, source Conjecture 24).} With the source's unresolved hypotheses and normalization, the assertion is
\[
\rho(X)=\operatorname{ord}_{t=q^{-1}}P_2(X,t),
\]
where \(\operatorname{ord}_{t=a}P(t)\) is the multiplicity of \(a\) as a root of a nonzero polynomial \(P\). The simultaneous use of geometric Frobenius, the Tate twist \(\Q_\ell(1)\), and the root \(q^{-1}\) requires source verification: the record does not establish a compatible normalization. The displayed equality preserves the supplied formulation without silently changing the twist, Frobenius convention, or root.
\subsection{Short English statement}
The supplied version of Tate's Conjecture T equates the Picard number of a surface \(X\) with the multiplicity of \(q^{-1}\) in a Frobenius polynomial. The hypotheses on the surface and base field are missing, and the polynomial convention, Tate twist, and stated root require reconciliation before a precise conjecture can be finalized.
\subsection{Sources}
\begin{itemize}
\item Supplied local UnsolvedMath registry record: permanent ID \(800712\), exact source problem code \texttt{catalogue entry}, AMR Open Problem Lists. The registry labels the problem open; that label is supplied metadata, not an independently verified literature assessment.
\item Auel, Brussel, Garibaldi, Vishne, \emph{Open Problems on Central Simple Algebras} (2010), Conjecture 24, identified in the record as Tate's Conjecture T. Supplied source URL: \url{https://arxiv.org/abs/1006.3304}. The record reports extraction from arXiv LaTeX and an access date of 2026-10-16; this is supplied provenance, not an independently performed source visit.
\item Tate, reference label \texttt{[Ta65]}, \(\S 3\), as cited in the supplied statement. No complete bibliographic entry for this reference is supplied.
\end{itemize}
\end{document}
